\section*{7. 中和と塩}
\begin{enumerate}
    \item 酸とアルカリの水溶液を混ぜたとき、たがいの性質を打ち消し合う反応を中和という。このとき、酸の陽イオンとアルカリの陰イオンが結びついて \textbf{水} ができる。この反応をイオン式を用いて表しなさい。\\[0.2cm]
    答（ \textcolor{red}{$\mathbf{H^+ + OH^- \rightarrow H_2O}$} ）
    \item 中和のとき、水のほかにできる、酸の陰イオンとアルカリの陽イオンが結びついた物質を何というか。\\[0.2cm]
    （ \textcolor{red}{\textbf{塩（えん）}} ）
    \item 次の酸とアルカリを混ぜて中和させたとき、起こる化学変化を \textbf{化学反応式} で表しなさい。
    \begin{enumerate}
        \item 硫酸と水酸化バリウム水溶液 \\[0.3cm]
        答（ \textcolor{red}{$\mathbf{H_2SO_4 + Ba(OH)_2 \rightarrow BaSO_4 + 2H_2O}$} ）\\[0.3cm]
        \item 塩酸と水酸化ナトリウム水溶液 \\[0.3cm]
        答（ \textcolor{red}{$\mathbf{HCl + NaOH \rightarrow NaCl + H_2O}$} ）
    \end{enumerate}
\end{enumerate}
